\section{Conclusion}
\label{sec:conclusion}

We compared a monolith with three decompositions of identical code, on a local cluster and on Amazon EKS, under an
equal CPU budget and a demand-proportional allocation and at two levels of service demand.
We also added autoscalers under a shared CPU ceiling in the cloud.
Decomposition added latency at low load: the p99 of the decomposed configurations was \xo{h1/min}--\xo{h1/max}
times the monolith's locally and \xo{eks/h1/min}--\xo{eks/h1/max} times in the cloud, mostly inter-service transit and, in the cloud, also longer database time.
It cut capacity to \xo{peakpct/min}--\xo{peakpct/max}\% (local) and \xo{eks/peakpct/min}--\xo{eks/peakpct/max}\%
(cloud) of the monolith's, because the edge service exhausted its CPU share while the rest of the budget stayed
idle. Allocating the budget by demand raised this to at most \xo{prop/peakpct/max}\%, and with 1\,ms of CPU in every
module operation to at most \xo{workprop/peakpct/max}\%. A sidecar mesh added
\xo{mesh/perhop/min}--\xo{eks/mesh/perhop/max}\,ms per call and more CPU than the application. With autoscalers the
monolith sustained \xo{hpa/sustained/browse/A} browse requests per second on one replica; the decomposed
configurations sustained every level only up to \xo{hpa/sustained/browse/B}, and B the monolith's level in one
repetition, at the ceiling. No decomposed configuration overtook the monolith
in any setting.

For a team deciding how finely to split, these measurements put a price on each synchronous call and show where it
is paid. They do not show a load at which decomposition pays for itself in performance; its case rests on the other
benefits it brings, and the price is smallest when services are coarse and do substantial work. The testbed, data
and analysis are published so that the same measurements can be repeated on other infrastructure.
